\centerline{\bf \Huge Инвариант}


\begin{flushright}
        \scriptsize{Эпизод 6. Схватка в Токио-3}
\end{flushright}

\problem{1}
На доске записаны все натуральные числа от 1 до 2022. Каждую минуту Игорь заменяет каждое из чисел на доске на сумму его цифр. Через несколько минут все числа стали однозначными. Каких чисел больше - пятёрочек или восьмёрочек?

\problem{2}
На хоккейной площадке лежат три шайбы, пронумерованные числами $1,2,3$. Хоккеист Ваня выбирает какую-то из них и бьет по ней так, что она пролетает между двумя другими. Он делает такие удары снова и снова. Может ли так оказаться, что после 2023 ударов каждая шайба лежит на своем первоначальном месте?

\problem{3}
В каждой клетке таблицы $n \times n$ стоит либо «+» либо «-». Разрешается выбрать произвольную клетку и поменять знаки в столбце и строке, содержащих ее (при этом знак в самой клетке тоже меняется). Найдите все $n$, при которых из любого изначального расположения знаков можно получить таблицу, заполненную плюсами.

\problem{4}
На доске записано несколько натуральных чисел. Каждую минуту Леша стирает с доски какие-то два числа и записывает вместо них их наибольший общий делитель и наименьшее общее кратное. Докажите, что когда-нибудь числа на доске перестанут меняться.

\problem{5}
В каждой вершине правильного 2022 -угольника стоит фишка. Разрешается за один ход сдвинуть любые две фишки в соседние с ними вершины. Можно ли с помощью таких операций собрать все фишки в одной вершине?

\problem{6}
Дно прямоугольной коробки было выложено плитками размерами $2 \times 2$ и $1 \times 4$. Плитки высыпали из коробки и при этом потеряли одну плитку $2 \times 2$. Вместо неё удалось достать плитку $1 \times 4$. Докажите, что теперь выложить дно коробки плитками не удастся.

\problem{7}
На прямой стоят две фишки: слева красная справа синяя. Разрешается производить любую из двух операций: вставку двух фишек одного цвета подряд (между фишками или с краю) и удаление пары соседних одноцветных фишек (между которыми нет других фишек). Можно ли с помощью таких операций оставить на прямой ровно две фишки: слева синюю, а справа красную?

\problem{8}
На доске $15 \times 15$ стоят 15 не бьющих друг друга ладей. Каждую ладью передвинули ходом коня 15 раз. Докажите, что теперь какие-то две ладьи бьют друг друга.

\newpage

\hspace{3mm}

\problem{9}
Четыре кузнечика сидят в вершинах квадрата. Каждую минуту один из них прыгает в точку, симметричную ему относительно другого кузнечика. Докажите, что кузнечики не могут в некоторый момент оказаться в вершинах квадрата большего размера.